\documentclass[a4paper,12pt]{article}

\usepackage[T2A]{fontenc}
\usepackage[utf8]{inputenc}
\usepackage[russian]{babel}

\usepackage{amsmath, amssymb, amsthm, mathtools}
\usepackage{helvet}
\renewcommand{\familydefault}{\sfdefault}
\usepackage{icomma}
\usepackage{geometry}
\usepackage{microtype}
\usepackage{tikz}
\usepackage{tkz-euclide}
\usepackage{pgfplots}
\pgfplotsset{compat=1.18}
\usepackage{tabularx, booktabs, longtable}
\usepackage{enumitem}
\usepackage{hyperref}
\usepackage{parskip}
\usepackage{fancyhdr}
\usepackage{setspace}
\usepackage{xcolor}

\definecolor{accent}{HTML}{008080}
\definecolor{darkaccent}{HTML}{004D4D}
\definecolor{textgray}{HTML}{333333}

\geometry{left=2.5cm, right=2.5cm, top=2.5cm, bottom=2.5cm}
\onehalfspacing

\begin{document}

% Титульный лист
\begin{titlepage}
\centering
\vspace*{2cm}
{\Large\bfseries\color{darkaccent} Чек-лист}\par
\vspace{0.5cm}
{\large Тема: Планиметрия. Треугольники и окружности}\par
\vspace{2cm}
{\large Лёвин Артём Александрович}\par
{\large эксклюзивно для Яны}\par
\vspace{2cm}
{\today}\par
\vfill
\begin{tikzpicture}
\draw[color=accent, line width=1.5pt] 
(0,0) .. controls (2,1.5) and (4,-0.5) .. (6,0.5) .. controls (8,1.5) and (10,-0.5) .. (12,0);
\end{tikzpicture}
\vspace{1cm}
\end{titlepage}

\section*{Чек-лист: Планиметрия}

\subsection*{Основные понятия}
\begin{itemize}[leftmargin=*]
\item \textbf{Синус угла} в прямоугольном треугольнике: $\sin A = \frac{\text{противолежащий катет}}{\text{гипотенуза}}$.
\item \textbf{Косинус угла}: $\cos A = \frac{\text{прилежащий катет}}{\text{гипотенуза}}$.
\item \textbf{Теорема синусов}: $\frac{a}{\sin A} = \frac{b}{\sin B} = \frac{c}{\sin C} = 2R$, где $R$ --- радиус описанной окружности.
\item \textbf{Теорема косинусов}: $c^2 = a^2 + b^2 - 2ab\cos C$.
\item \textbf{Вписанный угол} равен половине дуги, на которую он опирается.
\item \textbf{Центральный угол} равен градусной мере дуги, на которую он опирается.
\end{itemize}

\subsection*{Ключевые теоремы}
\begin{enumerate}[leftmargin=*]
\item \textbf{Теорема Пифагора}: $c^2 = a^2 + b^2$ (для прямоугольного треугольника).
\item \textbf{Свойство вписанного четырёхугольника}: сумма противоположных углов равна $180^\circ$.
\item \textbf{Признак касательной}: касательная перпендикулярна радиусу, проведённому в точку касания.
\end{enumerate}

\subsection*{Методы решения}
\begin{enumerate}[leftmargin=*]
\item Для нахождения сторон через углы использовать определения синуса и косинуса.
\item Если известны отношение сторон и одна сторона --- вводить коэффициент пропорциональности $x$.
\item Для связи сторон и углов применять теоремы синусов и косинусов.
\item При работе с окружностями использовать свойства вписанных и центральных углов.
\item Проверять ОДЗ: основание логарифма $a > 0$, $a \neq 1$; подкоренное выражение $\geq 0$.
\end{enumerate}

\subsection*{Типичные ошибки}
\begin{itemize}[leftmargin=*]
\item[$\times$] Путают противолежащий и прилежащий катеты.
\item[$\times$] Забывают, что синусу $\alpha$ соответствуют два угла: $\alpha$ и $180^\circ - \alpha$.
\item[$\times$] Не учитывают, что в геометрических задачах длины положительны.
\item[$\times$] При переходе через кратный корень знак производной не меняется.
\item[$\times$] Не проверяют принадлежность ответа заданному интервалу.
\end{itemize}

\newpage
\section*{Тренировочный блок}

\textbf{Уровень 1 (базовый)}

1. В треугольнике $ABC$ угол $C = 90^\circ$, $AB = 10$, $\sin A = 0{,}6$. Найдите катет $BC$.

2. Найдите $\cos A$, если в прямоугольном треугольнике $AC = 8$, $AB = 10$.

3. Градуcная мера дуги $AB$ равна $100^\circ$. Найдите вписанный угол, опирающийся на эту дугу.

\medskip
\textbf{Уровень 2 (средний)}

4. В треугольнике $ABC$: $BC = 3$, $\cos A = \frac{2\sqrt{5}}{5}$. Найдите $AC$.

5. Радиус описанной окружности равен $\sqrt{6}$, хорда $AC = 3$. Найдите тупой вписанный угол, опирающийся на эту хорду.

6. Найдите абсциссу точки, в которой касательная к графику функции параллельна оси $Ox$, если график производной пересекает ось $x$ в точке $x = -7$.

\medskip
\textbf{Уровень 3 (выше среднего)}

7. В трапеции $ABCD$ с основаниями $BC$ и $AD$ точки $E$ и $F$ --- середины сторон $BC$ и $AD$. В четырёхугольники $ABEF$ и $ECDF$ можно вписать окружности. Докажите, что трапеция равнобедренная.

8. Стороны треугольника относятся как $2:\sqrt{5}$. Периметр равен 20. Найдите стороны треугольника.

9. В треугольнике $ABC$ проведена биссектриса $AD$. Известно, что $AB = 6$, $AC = 9$, $BC = 10$. Найдите $BD$ и $DC$.

\medskip
\textbf{Уровень 4 (сложный)}

10. В равнобедренной трапеции с боковой стороной 7 в четырёхугольник, образованный серединами оснований и вершинами, можно вписать окружность радиуса 2{,}5. Найдите радиус окружности, описанной около трапеции.

\newpage
\section*{Ответы}

\begin{center}
\renewcommand{\arraystretch}{1.5}
\begin{tabular}{|c|l|}
\toprule
№ & Ответ \\
\midrule
1 & $BC = 6$ \\
2 & $\cos A = 0{,}8$ \\
3 & $50^\circ$ \\
4 & $AC = 6$ \\
5 & $120^\circ$ \\
6 & $x = -7$ \\
7 $AB = CD$ (суммы противоположных сторон равны) \\
8 & $4;\ 2\sqrt{5};\ 2\sqrt{5}$ \\
9 & $BD = 4,\ DC = 6$ \\
10 & $R = \frac{49}{12}$ \\
\bottomrule
\end{tabular}
\end{center}

\end{document}
